\section{Related Work}
\label{sec:related}
\subsection{Image Synthesis}\label{sec:related_synthesis}
深度生成模型~\cite{KingmaW14, Vahdat2020nvae, goodfellow2014generative, SAGAN, Song19generative,dhariwal2021diffusion,Oord16pixelcnn, Parmar18imagetransformer}在图像合成任务中取得了许多成功。基于GAN的方法在生成高保真样本方面展现了惊人的能力~\cite{goodfellow2014generative, biggan, Karras2019stylegan2, SAGAN, tseng2021regularizing}。
相比之下，基于似然的方法，如变分自编码器（VAE）~\cite{KingmaW14, Vahdat2020nvae}、扩散模型~\cite{Song19generative,dhariwal2021diffusion,ho2021cascaded}和自回归模型~\cite{Oord16pixelcnn, Parmar18imagetransformer}，提供了分布覆盖，因此可以生成更多样化的样本~\cite{Song19generative, Vahdat2020nvae, Oord16pixelcnn}。 

然而，直接在像素空间中最大化似然可能很有挑战性。因此，VQVAE~\cite{Oord17vqvae, Razavi19vqvae2}提议在两个阶段内在潜在空间中生成图像。在第一阶段，称为\textbf{tokenization}（分词），它试图将图像压缩到离散潜在空间中，主要包含三个分量：\vspace{-2.5mm}
\begin{itemize}
    \item 一个编码器$E$，学习将图像$x \in \mathbb{R}^{H\times W \times 3}$分词为潜在嵌入$E(x)$，\vspace{-2.5mm}
    \item 一个码本$\mathbf{e}_k \in \mathbb{R}^D, k\in 1, 2, \cdots, K$，用于最近邻查找以将嵌入量化为视觉token，以及\vspace{-2.5mm}
    \item 一个解码器$G$，从视觉token $\mathbf{e}$预测重建图像$\hat{x}$。
\end{itemize}\vspace{-1mm}
%  The training objective of the first stage is 
% \begin{equation} \label{eq:1}
% \begin{split}
%  \mathcal{L}_{VQ} = {} & ||sg[E(x)] -\mathbf{e}||_2^2 + \beta ||sg[\mathbf{e}] - E(x)||_2^2 \\
%   &+ \mathcal{L}_{rec}(x,\hat{x}).
% \end{split}
% \end{equation}
% Here the terms represent the codebook loss penalizing token $\mathbf{e}$ far from embedding $E(x)$, the commitment loss constraining $E$ to stick to one token in the codebook, and the reconstruction loss between input and reconstruction. $sg$ is the stop-gradient operator and $\beta$ is a commitment loss hyper-parameter. VQVAE first introduces this training objective and uses $L_2$ for the reconstruction loss. Later VQGAN proposes the additions of the perceptual loss and the adversarial loss to achieve higher compression efficiency.

\noindent 在第二阶段，它首先使用深度自回归模型预测视觉token的潜在先验，然后使用第一阶段的解码器将token序列映射到图像像素。由于两阶段方法的有效性，许多方法都遵循了这个范式。DALL-E~\cite{Ramesh21dalle}使用Transformer~\cite{Vaswani17attention}来改进第二阶段的token预测。VQGAN~\cite{Esser21vqgan}在第一阶段添加了对抗损失和感知损失~\cite{johnson2016perceptual, zhang2018unreasonable}来改进图像保真度。与我们同期的工作VIM~\cite{vim2021}提议使用VIT骨干~\cite{dosovitskiy2021vit}来进一步改进分词阶段。由于这些方法仍然采用自回归模型，第二阶段的解码时间随token序列长度而变化。 

\subsection{具有双向Transformer的掩盖建模}
Transformer架构~\cite{Vaswani17attention}最初在自然语言处理中被提出，最近已扩展到计算机视觉~\cite{dosovitskiy2021vit, caron2021dino}。
Transformer由多个自注意力层组成，允许捕获序列中所有元素对之间的交互。
特别是，BERT~\cite{Devlin19bert}引入了掩盖语言建模（MLM）任务用于语言表示学习。BERT~\cite{Devlin19bert}中使用的双向自注意力允许使用来自两个方向的上下文来预测MLM中的掩盖token。
在视觉中，BERT~\cite{Devlin19bert}中的掩盖建模已扩展到图像表示学习~\cite{he2021mae, Bao2022Beit}，其中图像被量化为离散token。然而，很少有工作成功地将相同的掩盖建模应用于图像生成~\cite{zhang2021ufcbert}，因为使用双向注意进行自回归解码很困难。
据我们所知，本文提供了第一个证据，证明了掩盖建模对于在常见ImageNet基准上的图像生成的有效性。
我们的工作受到自然语言处理中的双向机器翻译~\cite{ghazvininejad2019maskpredict, gu2020fully, gu2017non}的启发，我们的新颖性在于提议的新掩盖策略和解码算法，如我们的实验所证实的，这些对图像生成至关重要。
